\documentclass[10pt,aspectratio=169]{beamer}

\usetheme{metropolis}
\usepackage[utf8]{inputenc}
\usepackage[T1]{fontenc}
\usepackage[english]{babel}
\usepackage{listings}
\usepackage{booktabs}
\usepackage{amsmath,amssymb}

\definecolor{gfblue}{HTML}{1F4E79}
\definecolor{codebg}{HTML}{F4F4F2}
\definecolor{warn}{HTML}{B03A2E}
\definecolor{good}{HTML}{1E6B3A}

\setbeamercolor{frametitle}{bg=gfblue}
\setbeamercolor{alerted text}{fg=warn}

\lstset{
  basicstyle=\ttfamily\scriptsize,
  backgroundcolor=\color{codebg},
  frame=none,
  columns=fullflexible,
  keepspaces=true,
  showstringspaces=false,
  breaklines=true,
  xleftmargin=2mm,
  literate=
    {á}{{\'a}}1 {é}{{\'e}}1 {í}{{\'i}}1 {ý}{{\'y}}1 {ú}{{\'u}}1
    {č}{{\v{c}}}1 {ě}{{\v{e}}}1 {ň}{{\v{n}}}1 {ř}{{\v{r}}}1
    {š}{{\v{s}}}1 {ž}{{\v{z}}}1 {ů}{{\r{u}}}1
    {å}{{\aa}}1 {ä}{{\"a}}1 {ö}{{\"o}}1 {ü}{{\"u}}1 {Ä}{{\"A}}1 {Ö}{{\"O}}1
    {è}{{\`e}}1 {ê}{{\^e}}1
}

\newcommand{\bad}[1]{\textcolor{warn}{#1}}
\newcommand{\ok}[1]{\textcolor{good}{#1}}

\title{Porting Informath to a New Language}
\subtitle{GF-RGL and Claude: a Finnish case study}
\author{Aarne Ranta}
\date{July 2026}

\begin{document}

\maketitle

\begin{frame}{Outline}
  \tableofcontents
\end{frame}

% =====================================================================
\section{What Informath is}

\begin{frame}{The problem}
  Mathematics exists in two registers.

  \begin{columns}[T]
    \begin{column}{0.48\textwidth}
      \textbf{Formal}
      \begin{itemize}
        \item Agda, Rocq, Lean, Dedukti
        \item machine-checkable
        \item unreadable to most humans
      \end{itemize}
    \end{column}
    \begin{column}{0.48\textwidth}
      \textbf{Informal}
      \begin{itemize}
        \item English, French, German, Swedish
        \item readable
        \item ambiguous, unverifiable
      \end{itemize}
    \end{column}
  \end{columns}

  \vspace{1em}
  Informath translates between them --- in \emph{all} directions.
\end{frame}

\begin{frame}{Four directions}
  \begin{itemize}
    \item formal $\rightarrow$ informal \quad (\textbf{informalization})
    \item informal $\rightarrow$ formal \quad (\textbf{autoformalization})
    \item informal $\rightarrow$ informal \quad (translation, \emph{via} the formal)
    \item formal $\rightarrow$ formal \quad (in special cases)
  \end{itemize}

  \vspace{1em}
  Two goals:
  \begin{enumerate}
    \item complete informalization of anything expressible in Dedukti;
    \item generate \emph{and analyse} informal mathematical language,
          in all of its variation.
  \end{enumerate}

  \vspace{0.5em}
  \small Inspired by Ganesalingam, \emph{The Language of Mathematics} (2013).
\end{frame}

\begin{frame}[fragile]{One statement, many languages}
\begin{lstlisting}
Dedukti: prop110 : (a : Elem Int) -> (c : Elem Int) ->
  Proof (and (odd a) (odd c)) ->
  Proof (forall Int (b => even (plus (times a b) (times b c)))).
\end{lstlisting}

\vspace{-0.2em}
\begin{itemize}\footnotesize
  \item \textbf{En}: Let $a$ and $c$ be integers. Assume that $a$ and $c$ are odd.
        Then $ab + bc$ is even for all integers $b$.
  \item \textbf{De}: Seien $a$ und $c$ ganze Zahlen. Nimm an, dass $a$ und $c$ ungerade sind.
        Dann ist $ab + bc$ gerade für jede ganze Zahl $b$.
  \item \textbf{Sv}: Låt $a$ och $c$ vara heltal. Anta att $a$ och $c$ är udda.
        Då är $ab + bc$ jämnt för alla heltal $b$.
  \item \alert{\textbf{Fi}: \emph{--- what this talk adds}}
  \end{itemize}

  \vspace{0.2em}
  \footnotesize All generated from the Dedukti source. Agda, Rocq and Lean too.
\end{frame}

\begin{frame}{Architecture}
  \begin{center}
  \begin{tabular}{c}
    \textbf{Dedukti} --- the semantic core \\
    $\updownarrow$ \\
    \textbf{MathCore} --- an abstract syntax close to the logic \\
    $\updownarrow$ \\
    \textbf{Informath} --- adds variation, synonyms, symbolic/verbal choices \\
    $\updownarrow$ \\
    \textbf{Eng \quad Fre \quad Ger \quad Swe \quad \alert{Fin}} \\
  \end{tabular}
  \end{center}

  \vspace{1em}
  Agda, Rocq and Lean hang off Dedukti by (partial) conversions.

  The natural languages hang off \textbf{Informath} by GF concrete syntaxes.
\end{frame}

% =====================================================================
\section{GF in five minutes}

\begin{frame}[fragile]{Abstract vs.\ concrete syntax}
  GF separates \emph{what can be said} from \emph{how it is said}.

\begin{lstlisting}
abstract MathCore = {
  cat Prop ; Kind ; Ident ;
  fun CoreAllProp : Kind -> Ident -> Prop -> Prop ;
}
\end{lstlisting}

\begin{lstlisting}
concrete MathCoreEng of MathCore = {
  lincat Prop = S ;
  lin CoreAllProp kind x prop = ... "for all" ... ;
}
\end{lstlisting}

  \vspace{0.5em}
  One abstract syntax, $n$ concrete syntaxes $\Rightarrow$ $n^2$ translations,
  all going through the meaning.
\end{frame}

\begin{frame}[fragile]{The RGL, and functors}
  The Resource Grammar Library gives you, per language:
  \begin{itemize}
    \item a \textbf{Syntax API} --- \texttt{mkCN}, \texttt{mkNP}, \texttt{mkS} --- same names everywhere
    \item a \textbf{Paradigms API} --- \texttt{mkN}, \texttt{mkA}, \texttt{mkV} --- language-specific
  \end{itemize}

  So you write the linearization \emph{once}, against interfaces:

\begin{lstlisting}
incomplete concrete MathCoreFunctor of MathCore = open Syntax, Utilities in { ... }

concrete MathCoreFin of MathCore = MathCoreFunctor with
  (Syntax = SyntaxFin), (Utilities = UtilitiesFin) ;
\end{lstlisting}

  \vspace{0.3em}
  ``Functors make it maximally easy to port GF grammars into new languages.''
  \hfill --- \emph{Informath under the Hood}
\end{frame}

\begin{frame}{What a new language needs}
  \small
  \begin{center}
  \begin{tabular}{ll}
    \toprule
    \texttt{UtilitiesFin} & function words, \emph{instance} of the \texttt{Utilities} interface \\
    \texttt{CategoriesFin}, \texttt{MathCoreFin}, \dots & functor instantiations (one-liners) \\
    \texttt{VerbalConstantsFin} & $\sim$120 mathematical words \\
    \texttt{WikidataWordsFin} & $\sim$2000 mathematical words \\
    \texttt{InformathFin} & glue \\
    \bottomrule
  \end{tabular}
  \end{center}

  \vspace{1em}
  \begin{center}
  \Large Three one-liners and two dictionaries.
  \end{center}

  \vspace{0.5em}
  \begin{center}
  For Finnish, this is \ok{actually true}.
  \end{center}
\end{frame}

% =====================================================================
\section{Finnish: the RGL is not the problem}

\begin{frame}{The Finnish RGL is mature}
  \small
  \begin{center}
  \begin{tabular}{lrr}
    \toprule
    & \textbf{Finnish} & \textbf{Czech} \\
    \midrule
    \texttt{Paradigms} constructors & 55 & 7 \\
    \texttt{Paradigms} lines        & 1090 & 138 \\
    \texttt{Sentence} lines         & 74 & 32 \\
    \texttt{Idiom} lines            & 96 & 5 \\
    \texttt{Question} lines         & 99 & 7 \\
    \texttt{Structural} lines       & 334 & 31 \\
    \midrule
    \texttt{Missing<Lang>} stubs    & \ok{no such module} & \bad{97 \texttt{notYet}} \\
    \bottomrule
  \end{tabular}
  \end{center}

  \vspace{0.8em}
  Finnish has \texttt{mkV}, \texttt{mkVS}, \texttt{mkV3}, \texttt{mkPN},
  \texttt{mkSubj}, \texttt{mkPredet}, \dots{}

  Czech has \texttt{mkN}, \texttt{mkA}, \texttt{mkA2}, \texttt{mkV2},
  \texttt{mkAdv}, \texttt{mkPrep}, \texttt{mkConj}. \bad{No \texttt{mkV} at all.}

  \vspace{0.5em}
  \small ``Language $X$ is in the RGL'' is not a binary fact.
\end{frame}

\begin{frame}[fragile]{And it shows}
  Finnish morphology, for free, from a one-word paradigm:

\begin{lstlisting}
lin ring_Noun = mkNoun "rengas" ;
\end{lstlisting}

  \vspace{0.5em}
  \begin{center}
    \texttt{CoreAllProp (NounKind ring\_Noun) (StrIdent "R") ...} \\[0.4em]
    $\Longrightarrow$ \quad \ok{\emph{kaikille renkaille} $R$, \dots}
  \end{center}

  \vspace{0.8em}
  Allative plural, with consonant gradation \emph{ng} $\rightarrow$ \emph{nk}
  and the \emph{-as} stem. Nobody wrote that down.

  \vspace{0.5em}
  \small
  Compare Czech, where the one-argument \texttt{mkN} sends everything in
  \texttt{-ce} --- \emph{matice}, \emph{funkce}, \emph{kružnice} --- to the
  \emph{soudce} `judge' paradigm, and declines them as animate masculine.
\end{frame}

\begin{frame}{So the port should be easy}
  \begin{center}
  \Large It was. \\[1em]
  \normalsize And that is the whole problem.
  \end{center}
\end{frame}

% =====================================================================
\section{When the compiler stops testing you}

\begin{frame}[fragile]{Day one: it compiled}
  The Finnish modules were created as ``compilable templates'', by copying.

\begin{lstlisting}
$ gf InformathFin.gf
linking ... OK
\end{lstlisting}

  \vspace{0.5em}
  What was actually in them:

\begin{lstlisting}
-- UtilitiesFin.gf                    -- Swedish
then_Adv       = mkAdv "så" ;
such_that_Subj = mkSubj "så att" ;
arbitrary_A    = mkA "godtycklig" ;
iff_Subj       = mkSubj "om och endast om" ;
proof_Str      = "bevis" ;

-- VerbalConstantsFin.gf              -- English
lin natural_Noun = mkNoun "natural" "number" ;
lin set_Fam      = mkFam "set" ;
lin union_FunC   = mkFun "union" ;
\end{lstlisting}
\end{frame}

\begin{frame}{The central asymmetry}
  \begin{center}
  \begin{tabular}{ll}
    \toprule
    \textbf{Czech} & \textbf{Finnish} \\
    \midrule
    RGL half-built & RGL mature \\
    compiler screams & compiler silent \\
    $\sim$20 crash-driven iterations & \texttt{linking ... OK} from the start \\
    errors are \emph{structural} & errors are \emph{semantic} \\
    ``does it build?'' & ``is it Finnish?'' \\
    \bottomrule
  \end{tabular}
  \end{center}

  \vspace{1.2em}
  \begin{center}
  \Large A complete RGL removes your only automatic test.
  \end{center}

  \vspace{0.8em}
  \small
  In Czech, every mistake I made stopped the build. In Finnish, a grammar that
  emits fluent Swedish is indistinguishable, to the compiler, from a correct one.
\end{frame}

% =====================================================================
\section{Three layers of lexicon}

\begin{frame}{Layer 1: \texttt{UtilitiesFin} --- the function words}
  $\sim$23 strings. Swedish $\rightarrow$ Finnish, using the English \texttt{oper}
  identifiers as the only clue to meaning.

  \vspace{0.5em}
  \small
  \begin{center}
  \begin{tabular}{lll}
    \toprule
    identifier & was (sv) & now (fi) \\
    \midrule
    \texttt{then\_Adv}       & så & niin \\
    \texttt{such\_that\_Subj} & så att & siten että \\
    \texttt{arbitrary\_A}    & godtycklig & mielivaltainen \\
    \texttt{iff\_Subj}       & om och endast om & jos ja vain jos \\
    \texttt{proof\_Str}      & bevis & todistus \\
    \texttt{theorem\_Str}    & teorem & lause \\
    \texttt{definition\_Str} & definition & määritelmä \\
    \texttt{infinity\_NP}    & \bad{äärettön} \emph{(a typo, in Swedish's place)} & ääretön \\
    \bottomrule
  \end{tabular}
  \end{center}

  \vspace{0.5em}
  \footnotesize
  Some strings were \emph{already} Finnish (määritellä, olettaa, tyyppi, funktio,
  joukko). The file was a sediment of two languages, and compiled either way.
\end{frame}

\begin{frame}{Layer 2: \texttt{VerbalConstantsFin} --- the maths words}
  $\sim$120 strings, an English copy with two Finnish words in it.

  \vspace{0.5em}
  \begin{itemize}\small
    \item \textbf{Types}: propositio, numero, totuusarvo, luonnollinen luku,
          rationaaliluku, reaaliluku, joukko
    \item \textbf{Arithmetic}: summa, erotus, tulo, osamäärä,
          \emph{korotus $\dots$ potenssiin}, vastaluku,
          \emph{logaritmi $\dots$ kannassa}, neliöjuuri, itseisarvo, kertoma,
          suurin yhteinen tekijä
    \item \textbf{Sets}: yhdiste, leikkaus, komplementti, karteesinen tulo,
          potenssijoukko, osajoukko / aito osajoukko, ylijoukko, alkio, tyhjä joukko
    \item \textbf{Analysis}: pinta-ala, säde, kohtisuora, pituus, normi,
          mahtavuus, aste, sini / kosini / tangentti, pistetulo, kulma
  \end{itemize}

  \vspace{0.5em}
  \small
  Note \emph{luonnollinen luku}: the template had
  \texttt{mkNoun "natural" "number"} --- two arguments, adjective plus noun.
  The structure was right. Only the language was wrong.
\end{frame}

\begin{frame}[fragile]{Layer 3: \texttt{WikidataWordsFin} --- and the one thing that \emph{did} fail}
  2051 lins, byte-identical to the English module.

  It was the only file that did not compile --- and so it had simply been
  \textbf{commented out} of \texttt{InformathFin}:

\begin{lstlisting}
concrete InformathFin of Informath =
  MathCoreFin,
  VerbalConstantsFin,
----  WikidataWordsFin,        -- doesn't compile
  ProperNamesFin ;
\end{lstlisting}

  \vspace{0.5em}
  \textbf{Diagnosis.} It referenced the \emph{English morphodict}:

\begin{lstlisting}
lin 'knot_Noun'  = mkNoun knot_N ;        -- knot_N lives in ExtractEng
lin 'abelian_Adj' = mkAdj Abelian_A ;     -- MathWordsFin is a different, smaller set
\end{lstlisting}

  \vspace{0.3em}
  A grammar that ``worked'' by excluding 2000 of its 2100 words.
\end{frame}

\begin{frame}[fragile]{The fix: drop the morphodict}
  Finnish has good smart paradigms. So do not reference lemmas at all ---
  inline the strings:

\begin{lstlisting}
concrete WikidataWordsFin of WikidataWords = CategoriesFin **
  open UtilitiesFin in {          -- no ExtractFin

lin 'knot_Noun'        = mkNoun "solmu" ;
lin 'semigroup_Noun'   = mkNoun "puoliryhmä" ;
lin 'quasicoherent_Adj' = mkAdj "kvasikoherentti" ;
lin 'close_Verb'       = mkVerb "sulkea" ;
\end{lstlisting}

  \vspace{0.5em}
  \begin{itemize}
    \item \texttt{ParadigmsFin.mkN} handles gradation, vowel harmony, stem types.
    \item The module now needs no morphodict, and compiles standalone.
    \item \ok{Enabled in \texttt{InformathFin}.} \texttt{linking ... OK} --- this time meaning it.
  \end{itemize}

  \vspace{0.3em}
  \small This is exactly the move that is \bad{unsafe in Czech}, where the
  one-argument \texttt{mkN} guesses wrong for $\sim$19\% of nouns.
\end{frame}

\begin{frame}{Translating 2036 terms with an LLM}
  \begin{columns}[T]
  \begin{column}{0.52\textwidth}
  \textbf{Pipeline}
  \begin{enumerate}\small
    \item extract identifiers, split by category \\
          \small 1216 nouns, 749 adjectives, 86 verbs
    \item 7 batches, 7 subagents in parallel
    \item each returns \texttt{english\textbackslash tfinnish} TSV
    \item merge, check coverage
    \item code-generate the module
  \end{enumerate}
  \end{column}
  \begin{column}{0.44\textwidth}
  \textbf{Result}
  \begin{itemize}\small
    \item full coverage, 0 missing
    \item 0 malformed, 0 duplicate keys
    \item every abstract identifier preserved
  \end{itemize}
  \vspace{0.5em}
  \small Spot-checks: ryhmä, rengas, kunta, avaruus, monisto, lyhde (sheaf),
  ydin (kernel), ominaisarvo, kohomologia, jaoton (prime).
  \end{column}
  \end{columns}

  \vspace{1em}
  \begin{center}
  Nominative singular for nouns, A-infinitive for verbs. \\
  One column. Finnish morphology does the rest.
  \end{center}
\end{frame}

\begin{frame}[fragile]{It speaks Finnish}
\begin{lstlisting}
VarHypo (StrIdent "x") (NounKind group_Noun)
  ==>  olkoon $ x $ ryhmä

CoreAllProp (NounKind ring_Noun) (StrIdent "R") FalseProp
  ==>  kaikille renkaille $ R $ , meillä on ristiriita

CoreExistProp (NounKind matrix_Noun) (StrIdent "M") FalseProp
  ==>  on olemassa matriisi $ M $ , siten että meillä on ristiriita

CoreIfProp FalseProp FalseProp
  ==>  jos meillä on ristiriita , niin meillä on ristiriita

PropHypo FalseProp
  ==>  oleta , että meillä on ristiriita
\end{lstlisting}

  \vspace{0.3em}
  \ok{\emph{olkoon}} --- a real jussive, from \texttt{ImpP3} in \texttt{IdiomFin}. \\
  \small In Czech, \texttt{IdiomCze} was empty and \texttt{ImpP3} had to be written by hand.
\end{frame}

% =====================================================================
\section{What the compiler cannot see}

\begin{frame}[fragile]{Wrongness that links cleanly, I: prepositions}
  Finnish expresses these relations with \emph{cases}, not prepositions.

\begin{lstlisting}
applied_to_Prep : Prep = mkPrep "sovellettuna" ;
defined_as_Prep : Prep = mkPrep "määriteltynä" ;
as_Prep         : Prep = mkPrep "muodossa" ;
at_Prep         : Prep = mkPrep "kohdassa" ;
over_Prep       : Prep = mkPrep "yli" ;
\end{lstlisting}

  \vspace{0.5em}
  These are pre-positioned words that do not govern the case of their argument.

  \vspace{0.5em}
  They \textbf{compile}. They \textbf{read acceptably}. They are \textbf{not right}.

  \vspace{0.5em}
  \small
  No test in the repository distinguishes this from correct Finnish.
\end{frame}

\begin{frame}[fragile]{Wrongness that links cleanly, II: the lexicon}
  \begin{itemize}
    \item \textbf{Coinages.} $\sim$1500 terms have no established Finnish word.
      \begin{center}\ttfamily\small
      quiver $\longrightarrow$ nuolikko
      \end{center}
      Plausible. Invented. Unreviewed.
    \item \textbf{Capitalisation.} 8 proper-name adjectives kept a capital:
\begin{lstlisting}
lin Abelian_Adj = mkAdj "Abelin" ;    lin Baire_Adj = mkAdj "Bairen" ;
lin abelian_Adj = mkAdj "abelin" ;    -- the same word, twice, two ways
\end{lstlisting}
    \item \textbf{Irregular stems.} Inline strings rely on smart paradigms.
      A few will mis-inflect. Nothing reports which.
  \end{itemize}

  \vspace{0.5em}
  \begin{center}
  \alert{Every one of these is a silent defect in a grammar that links.}
  \end{center}
\end{frame}

\begin{frame}[fragile]{What copying hid}
  \texttt{WikidataWordsFin} was derived from \texttt{WikidataWordsEng}. So it
  inherited English's bookkeeping:

  \vspace{0.5em}
  \begin{center}
  \begin{tabular}{lrrrr}
    \toprule
    & lins & unique & \bad{not in abstract} & \bad{duplicated} \\
    \midrule
    \texttt{WikidataWordsEng} & 2051 & 2036 & 33 & 15 \\
    \texttt{WikidataWordsFin} & 2051 & 2036 & 33 & 15 \\
    \midrule
    \texttt{WikidataWordsCze} & \ok{2003} & \ok{2003} & \ok{0} & \ok{0} \\
    \bottomrule
  \end{tabular}
  \end{center}

  \vspace{0.3em}
  The 33 orphans are exactly the symbol tokens:

\begin{lstlisting}
lin A_Noun = mkNoun "A" ;   lin C222_Noun = mkNoun "C222" ;   lin csgn_Noun = mkNoun "csgn" ;
\end{lstlisting}

  \vspace{0.2em}
  \footnotesize
  Seven subagents dutifully translated identifiers that the abstract syntax does
  not contain. GF only \emph{warns}.
  The Czech module was generated \emph{from the abstract} --- so it cannot drift.
\end{frame}

% =====================================================================
\section{Lessons}

\begin{frame}{What the LLM was good at}
  \begin{itemize}
    \item \textbf{Bulk lexicography.} 2036 terms, 7 agents, minutes.
          The mathematical vocabulary is genuinely good: \emph{lyhde} for sheaf,
          \emph{ydin} for kernel, \emph{jaoton} for prime.
    \item \textbf{Diagnosis.} Working out that \texttt{WikidataWordsFin} failed
          because \texttt{knot\_N} lives in \texttt{ExtractEng} and
          \texttt{MathWordsFin} is a different, smaller extraction ---
          then proposing to drop the morphodict entirely.
    \item \textbf{Noticing the sediment.} That \texttt{UtilitiesFin} was
          Swedish-with-Finnish-patches, and that \texttt{äärettön} was a typo
          for \texttt{ääretön}.
    \item \textbf{Mechanical faithfulness.} Every abstract identifier preserved
          exactly, through a code-generated module.
  \end{itemize}
\end{frame}

\begin{frame}{What it was bad at}
  \begin{itemize}
    \item \textbf{Treating \texttt{linking ... OK} as success.} \\
          It is, at best, evidence that the file is well-formed GF.
    \item \textbf{Volunteering that \texttt{mkPrep "yli"} is not how Finnish works.} \\
          It compiled, so it was reported as done. The caveat exists because
          somebody asked afterwards.
    \item \textbf{Questioning its input.} It translated all 2051 lins, including
          33 that no \texttt{fun} references and 15 that were already there.
          It was given a file and it translated the file.
    \item \textbf{Knowing which coinages are real.} \emph{nuolikko} and
          \emph{lyhde} are presented with identical confidence. One is standard.
  \end{itemize}
\end{frame}

\begin{frame}{Lessons for the next language}
  \begin{enumerate}
    \item \textbf{Check the RGL first --- but for the opposite reason.} \\
          \small A \emph{poor} RGL costs you weeks of resource-grammar work.
          A \emph{good} RGL costs you your safety net.
    \item \normalsize \textbf{Generate from the abstract syntax, never copy a sibling.} \\
          \small Copying gave Finnish 33 dead lins and 15 duplicates, for free,
          silently. Code-generating Czech from \texttt{WikidataWords.gf} made
          ``complete'' a checkable property.
    \item \normalsize \textbf{Build the test the compiler will not give you.} \\
          \small Coverage, category, morphology class. Reject, do not repair.
    \item \normalsize \textbf{Smart paradigms are a property of the RGL, not of the LLM.} \\
          \small Trust them in Finnish. Do not trust them in Czech.
    \item \normalsize \textbf{Separate ``the model produced it'' from ``it is right''.} \\
          \small The 2036 Finnish terms have had no native review. That sentence
          belongs in the commit message.
  \end{enumerate}
\end{frame}

\begin{frame}{Still open}
  \begin{itemize}
    \item \textbf{Native review of $\sim$2036 translations}, especially the rare coinages.
    \item Convert the \texttt{mkPrep} string relations to case-governed prepositions.
    \item Lowercase the proper-name adjectives (\emph{Abelin} / \emph{abelin}), and
          decide which of the two duplicates survives.
    \item Add explicit paradigms where the smart guess mis-inflects an irregular stem.
    \item Purge the 33 orphan lins --- in English as well as Finnish.
  \end{itemize}
\end{frame}

\begin{frame}{Conclusion}
  \begin{center}
  \Large
  Czech was hard because the compiler kept stopping me.

  \vspace{0.8em}
  Finnish was hard because it never did.
  \end{center}

  \vspace{1.5em}
  \normalsize
  A mature RGL delivers exactly what it promises: the functor instantiates, the
  morphology is free, and a two-thousand-word lexicon is an afternoon.

  \vspace{0.5em}
  What it also delivers is a grammar that compiles, links, and speaks Swedish ---
  and no way to tell, except by reading it.

  \vspace{1em}
  \begin{center}
  \small
  \texttt{github.com/GrammaticalFramework/informath}
  \end{center}
\end{frame}

\end{document}
